\documentclass[11pt]{article}
\usepackage[T1]{fontenc}
\usepackage[utf8]{inputenc}
\usepackage{lmodern,amsmath,amssymb,amsthm,microtype}
\usepackage[margin=25mm]{geometry}
\usepackage{hyperref}
\hypersetup{hidelinks,pdftitle={A finite-energy curvature operator on curves over arbitrary complete fields}}
\setlength{\emergencystretch}{3em}
\newtheorem{theorem}{Theorem}
\newtheorem{lemma}[theorem]{Lemma}
\newtheorem{proposition}[theorem]{Proposition}
\newtheorem{corollary}[theorem]{Corollary}
\newcommand{\PSH}{\operatorname{PSH}}
\newcommand{\MA}{\operatorname{MA}}
\newcommand{\ddc}{dd^c}
\newcommand{\an}{\mathrm{an}}
\newcommand{\E}{\mathcal E}
\newcommand{\DF}{\mathcal D}
\title{A finite-energy curvature operator on curves\\over arbitrary complete fields}
\author{P16.4 research continuation 216\\Unreviewed partial result}
\date{7 October 2026}
\begin{document}
\maketitle
\begin{abstract}
For a smooth geometrically integral projective curve over a complete
characteristic-zero non-Archimedean field, we propose a direct extension
of the genuine continuous curvature measure to actual finite-energy
psh metrics. Trivial and nondiscrete valuations are included. The
argument controls ordered approximation nets by a Dirichlet seminorm,
uses dense continuous PL tests to obtain positive Radon measures, and
then proves integrability and weighted decreasing continuity through
an exact pairing formula. The potential remains an actual all-point
decreasing-FS-net endpoint; a Hilbert completion is used only for
estimates. No envelope continuity, general-field equation theorem or
higher-dimensional mixed operator is asserted. The entire new proof
awaits independent verification.
\end{abstract}

\section{Actual potentials, raw mass and the proposed theorem}

Let $K$ be any complete non-Archimedean real-valued field of
characteristic zero, with trivial, discrete or nondiscrete absolute
value. Let $X/K$ be a smooth geometrically integral projective curve,
and $L$ an actual ample rational line bundle, of degree $V>0$.
Put $T=X^{\an}$, a compact Hausdorff space, and choose a continuous
psh metric $\psi$ on $L$. The term psh has precisely
\cite[Definition 7.2]{BE}: a decreasing-net endpoint of continuous
Fubini--Study metrics on the actual fixed rational bundle. We exclude
the identically $-\infty$ metric. Write an absolute metric as $\psi+u$;
the scalar $u$ is evaluated at every point of $T$, including supported
points. A pole means its literal value $-\infty$.

For continuous psh $u$, let $\mu_{u,c}$ be the genuine curvature
measure of $\psi+u$ from \cite[Theorem 8.4]{BE}, and put
$\mu_\psi=\mu_{0,c}$. These are positive Radon measures of raw mass
$V$. Define the continuous relative energy and its extension by
\begin{equation}\label{eq:def}
 \begin{aligned}
 E(u)&=\frac12\int_Tu\,(\mu_{u,c}+\mu_\psi),\\
 E(u)&=\inf\{E(v):v\text{ continuous psh scalar},\ v\geq u\}
                \quad\text{for singular psh }u,\\
 \E^1&=\{u\text{ actual psh scalar}:E(u)>-\infty\}.
 \end{aligned}
\end{equation}
Every psh endpoint has a continuous majorant. Thus its energy is
defined in $[-\infty,\infty)$. The reference and the raw normalization
are fixed throughout. No synthetic element is substituted for a
potential in this definition.

\begin{theorem}\label{thm:main}
There is a unique positive Radon measure $\mu_u$ of mass $V$ for
each $u\in\E^1$ which is the weak limit of
$\mu_{u_i,c}$ for every decreasing continuous psh net $u_i\downarrow u$.
The limit is independent of that net and agrees with the genuine
continuous measure for continuous $u$.

Every $w\in\E^1$ is integrable against every $\mu_u$, and
\[
 E(u)=\frac12\int_Tu\,(\mu_u+\mu_\psi).
\]
If $u_i\downarrow u$ in $\E^1$ with $u\in\E^1$, then
$\mu_{u_i}\to\mu_u$ weakly. If also $w_i\downarrow w$ in $\E^1$
with $w\in\E^1$, then, on the common directed index set,
\[
 \int_T w_i\,d\mu_{u_i}\longrightarrow\int_Tw\,d\mu_u.
\]
The measures and all these statements are independent of the chosen
continuous psh reference, and scalar extension satisfies the ordinary
pushforward law with unchanged degree.

For actual $u,v\in\E^1$ and $u_t=(1-t)u+tv$, $0\leq t\leq1$,
the endpoints are finite-energy psh, their measures satisfy
$\mu_{u_t}=(1-t)\mu_u+t\mu_v$, and
\[
 \frac d{dt}E(u_t)=\int_T(v-u)\,d\mu_{u_t}.
\]
The integral is taken on the common finite-value full-measure set;
no subtraction of two infinities is used.
\end{theorem}

This is a curve operator theorem. In dimension one its only curvature
slot is the genuine degree-one Monge--Ampere operator. It provides
no full-original P16.4 equation existence theorem and no operator in
arbitrary higher dimension.

\section{The continuous Dirichlet form}

Let $\DF$ be the real vector space spanned by scalar differences of
continuous psh metrics on the same semiample rational line bundle,
allowing different bundles for different differences, together with
the constants. For $f,g\in\DF$, a chosen finite representation gives
curvature differences which are finite signed Radon measures. Compute
\[
 D(f,g)=-\int_T f\,\ddc g
\]
using such representations; independence is checked below.
All expressions are functions on the full carrier.
Finite linear combinations have bounded total variation of their
curvature measures.

\begin{lemma}\label{lem:dirichlet}
The curvature differences and $D$ are representation-independent.
The form $D$ is symmetric and positive semidefinite, with constants
in its kernel, and satisfies Cauchy--Schwarz. The space $\DF$ is
uniformly dense in $C(T)$. For continuous fixed-$L$ psh scalars,
\begin{equation}\label{eq:quadratic}
 E(u)=\int_Tu\,d\mu_\psi-\tfrac12D(u,u).
\end{equation}
If $u\leq v$ are such scalars, then
\begin{equation}\label{eq:ordered}
 0\leq D(v-u,v-u)\leq 2(E(v)-E(u)).
\end{equation}
\end{lemma}
\begin{proof}
Approximate each continuous psh metric uniformly by smooth
semipositive metrics, using \cite[Theorem 7.14]{BE}. For the resulting
smooth scalar differences, the compact Stokes rule gives
\[
 -\int_T f\,\ddc g=\int_Tdf\wedge d^cg,
 \qquad D(f,f)\geq0.
\]
These are the smooth differential and positivity conventions behind
\cite[Proposition 8.3]{BE}. Symmetry is its integration-by-parts
law. Uniform regularization and uniform-to-weak curvature convergence
pass the integrals to the continuous differences. Indeed for
$f_j\to f$ uniformly and curvature differences of $g_j$ weakly
convergent with a uniform total-variation bound, split the error into
$\int(f_j-f)\ddc g_j$ and $\int f(\ddc g_j-\ddc g)$.
The first tends to zero by the variation bound and the second by
weak convergence. Apply this simultaneously to every finite sum in
$\DF$. Positivity and symmetry persist, and the nonnegative
quadratic $D(f+tg,f+tg)$ proves Cauchy--Schwarz. Constant curvature
and integration by parts put constants in the kernel.

Every $\mathrm{PL}_{\mathbb R}$ function is a difference of FS
metrics on one common semiample rational bundle, by Definition 5.15
of \cite{BE}. These functions belong to $\DF$. Theorem 5.18, with
the nontrivial subgroup $\mathbb R$ even when $K$ is trivially valued,
proves uniform density. The complete separation argument in that
theorem includes supported seminorm points. For a representation
of the zero scalar with candidate curvature $\lambda$, the
integration-by-parts identity just proved gives
$\int h\,d\lambda=\int 0\,\ddc h=0$ for every $h\in\DF$.
Density and finite variation imply $\lambda=0$. Thus curvature
differences, and hence $D$, do not depend on representations.

Since $\mu_{u,c}=\mu_\psi+\ddc u$, the continuous energy definition
immediately gives \eqref{eq:quadratic}. Integration by parts gives
the exact difference formula
\[
 2(E(v)-E(u))=\int_T(v-u)(\mu_{v,c}+\mu_{u,c}).
\]
On the other hand
\[
 D(v-u,v-u)=\int_T(v-u)(\mu_{u,c}-\mu_{v,c}).
\]
When $v-u\geq0$, positivity of both curvature measures bounds the
second expression above by the first. This proves
\eqref{eq:ordered} without an envelope or comparison principle.
\end{proof}

Let $H$ be the Hilbert completion of $\DF$ modulo the null space of
$D$. We use $\|f\|_H^2=D(f,f)$. This construction is an estimate
space, not the definition of a psh potential or of its curvature
measure.

\section{Energy limits and Radon integration of directed nets}

\begin{lemma}\label{lem:energy}
If actual psh scalars $v_i\downarrow u$ and each $v_i$ has an
extended energy, then $E(v_i)\downarrow E(u)$.
In particular this holds for defining continuous psh nets and for
decreasing finite-energy nets. The endpoint may have energy
$-\infty$ in this lemma.
\end{lemma}
\begin{proof}
Monotonicity of extended energy follows directly from its majorant
infimum, and its constant-shift law follows from the corresponding
continuous law. Hence $E(v_i)\geq E(u)$. Fix any continuous psh
majorant $a\geq u$ and $\epsilon>0$. Each $v_i$ is upper
semicontinuous, so the sets
$\{x:v_i(x)-a(x)\geq\epsilon\}$ are compact. They decrease with
empty intersection. Compactness and directedness give
$v_i\leq a+\epsilon$ on a tail. Consequently
$\inf_iE(v_i)\leq E(a)+V\epsilon$. Let
$\epsilon\downarrow0$, then take the infimum over $a$.
\end{proof}

\begin{lemma}\label{lem:radon}
For a finite Radon measure $\lambda$ on a compact Hausdorff space,
an increasing directed net of nonnegative lower semicontinuous
functions $f_i$ satisfies
\[
 \int\sup_i f_i\,d\lambda=\sup_i\int f_i\,d\lambda.
\]
\end{lemma}
\begin{proof}
For $t\geq0$, the strict superlevel of $f=\sup_i f_i$ is
$\{f>t\}=\bigcup_i\{f_i>t\}$, a union of open sets.
Every compact subset of this union has a finite subcover, and one
directed index dominates that finite cover. Inner regularity gives
$\lambda(\{f>t\})=\sup_i\lambda(\{f_i>t\})$.
For a finite lower simple layer approximation to $\min(f,M)$,
choose one index simultaneously approximating its finitely many
strict-superlevel masses. Integrating the corresponding stage gives
the lower bound for that finite layer approximation. Refine the
finite levels and let $M\uparrow\infty$. The reverse bound follows
from $f_i\leq f$. This proves the identity, also for infinite
integrals. It is not a monotone-convergence rule for arbitrary
measurable nets.
\end{proof}

For any defining net $u_i\downarrow u$, a tail has a common finite
upper bound, since one continuous majorant is bounded on $T$.
Adding a constant to $\psi$ lets us assume that this whole tail and
its endpoint are nonpositive. Such a shift leaves $\mu_\psi$ and
$D$ unchanged and shifts the energy by that constant times $V$.
This harmless reduction will be used in estimates below.

\section{An actual endpoint has a canonical Dirichlet vector}

\begin{proposition}\label{prop:vector}
For every actual $u\in\E^1$, the class of any decreasing continuous
psh approximation $u_i$ converges strongly in $H$ to an element
denoted $[u]$, independent of the approximation. Moreover $u$ is
integrable against $\mu_\psi$ and
\begin{equation}\label{eq:singular-quadratic}
 E(u)=\int_Tu\,d\mu_\psi-\tfrac12\|[u]\|_H^2.
\end{equation}
\end{proposition}
\begin{proof}
Assume $u_i\leq0$ on a tail and put
$A(u_i)=-\int_Tu_i\,d\mu_\psi$ and $B(u_i)=-E(u_i)$.
Equation \eqref{eq:quadratic} reads
\[
 B(u_i)=A(u_i)+\tfrac12\|u_i\|_H^2,
 \qquad A(u_i)\geq0.
\]
Lemma \ref{lem:energy} gives $B(u_i)\uparrow B(u)<\infty$.
The increasing continuous functions $-u_i$ and Lemma
\ref{lem:radon} give $A(u_i)\uparrow A(u)\leq B(u)$.
Thus $u$ is integrable against the reference measure.

For two comparable indices $j\geq i$, \eqref{eq:ordered} gives
\[
 \|u_i-u_j\|_H^2\leq2(E(u_i)-E(u_j)).
\]
On a tail this is arbitrarily small, because the energies have a
finite limit. For two incomparable late indices pass to a common
later index and use the triangle inequality. The net is Cauchy in
$H$ and hence converges. Passing the displayed quadratic identity
to the limit proves \eqref{eq:singular-quadratic}.

For a second net $v_j\downarrow u$, use the product-directed net
$a_{i,j}=\max(u_i,v_j)$, continuous psh by Corollary 7.5 of
\cite{BE}. It decreases to $u$ at every point. Its energies tend
$E(u)$ by Lemma \ref{lem:energy}. Applying
\eqref{eq:ordered} to $u_i\leq a_{i,j}$ and
$v_j\leq a_{i,j}$ shows that both differences tend to zero in
$H$. Their Hilbert limits therefore coincide. The value of $u$ at
every point was fixed before this estimate; no identification of
all Hilbert elements with potentials is asserted.
\end{proof}

\section{A positive measure obtained from genuine curvature limits}

\begin{proposition}\label{prop:measure}
For $u\in\E^1$, the measures $\mu_{u_i,c}$ along every decreasing
continuous psh approximation converge weakly to one positive Radon
measure $\mu_u$ of mass $V$, characterized on $f\in\DF$ by
\begin{equation}\label{eq:tests}
 \int_Tf\,d\mu_u=\int_Tf\,d\mu_\psi-D(f,[u]).
\end{equation}
\end{proposition}
\begin{proof}
For each continuous approximation, integration by parts gives
$\int f\,d\mu_{u_i,c}=\int f\,d\mu_\psi-D(f,u_i)$.
Strong Hilbert convergence makes these integrals converge for every
$f\in\DF$. The limiting functional is positive on nonnegative
elements of $\DF$ and satisfies
$|\Lambda(f)|\leq V\|f\|_\infty$, since every approximating
curvature is positive of mass $V$. Uniform density extends it to
a bounded functional on $C(T)$. It is positive there as well:
approximate a nonnegative test uniformly by $f\in\DF$ and add
its uniform error as a constant to make $f$ nonnegative.
The functional takes $1$ to $V$. Riesz representation supplies a
positive Radon measure of mass $V$.

Density and the same mass bound show weak convergence for every
continuous test, not just $\DF$ tests. Independence of $[u]$ gives
independence of the measure. For a continuous endpoint, the
constant approximation or continuous uniform-to-weak law gives
its original genuine curvature. Uniqueness follows from the weak
limit requirement, or from the dense tests in \eqref{eq:tests}.
Thus the measure was constructed before choosing a target or
defining any unbounded weighted integral.
\end{proof}

\begin{proposition}\label{prop:weighted}
Every $w\in\E^1$ is integrable against every $\mu_u$, and
\begin{equation}\label{eq:weighted}
 \int_Tw\,d\mu_u=\int_Tw\,d\mu_\psi-D([w],[u]).
\end{equation}
For nonpositive $u,w$ relative to a common shifted reference,
\[
 0\leq\int_T(-w)\,d\mu_u
 \leq B(w)+2\sqrt{B(w)B(u)},\qquad B(v)=-E(v).
\]
\end{proposition}
\begin{proof}
Take continuous psh $w_j\downarrow w\leq0$. Each $w_j$ belongs
to $\DF$, so \eqref{eq:tests} gives
\[
 \int_T(-w_j)\,d\mu_u=A(w_j)+D(w_j,[u]).
\]
Lemma \ref{lem:radon} on this single Radon measure shows that
the left side increases to $\int_T(-w)\,d\mu_u$.
The right side tends to $A(w)+D([w],[u])$, a finite value by
Proposition \ref{prop:vector}. This proves integrability and
\eqref{eq:weighted}. The bounds $A(w)\leq B(w)$ and
$\|[v]\|_H^2\leq2B(v)$ give the asserted estimate by
Cauchy--Schwarz. Adding constants back proves the unrestricted
statement. Taking $w=u$ and using
\eqref{eq:singular-quadratic} gives the asserted genuine weighted
energy identity.
\end{proof}

\section{Decreasing singular nets and unbounded first variation}

\begin{lemma}\label{lem:ordered-singular}
For $u,v\in\E^1$ with $u\leq v$,
\[
 \|[v]-[u]\|_H^2\leq2(E(v)-E(u)).
\]
\end{lemma}
\begin{proof}
Choose decreasing continuous approximations $u_i\downarrow u$ and
$v_j\downarrow v$. The product-directed net
$a_{i,j}=\max(u_i,v_j)$ decreases to $\max(u,v)=v$ at every point.
Apply \eqref{eq:ordered} to $u_i\leq a_{i,j}$.
Propositions \ref{prop:vector} and Lemma \ref{lem:energy} pass
its two sides to the stated inequality.
\end{proof}

For an arbitrary decreasing finite-energy net $u_i\downarrow u$
with finite-energy endpoint, Lemma \ref{lem:energy} and
Lemma \ref{lem:ordered-singular} show $[u_i]\to[u]$ strongly in
$H$. The same common finite upper bound makes $-u_i$ increasing
nonnegative lower semicontinuous on a tail. Lemma \ref{lem:radon}
gives $\int u_i\,d\mu_\psi\to\int u\,d\mu_\psi$.
The dense test formula \eqref{eq:tests} then gives weak convergence
$\mu_{u_i}\to\mu_u$.

If simultaneously $w_i\downarrow w$ has a finite-energy endpoint,
its Hilbert vector and reference integral converge too. In
\eqref{eq:weighted}, the pairing
$D([w_i],[u_i])$ tends to $D([w],[u])$ by strong convergence
and bounded norms on tails. This proves the joint weighted limit
for singular nets. It does not use a measurable-net convergence
theorem against moving measures.

For two finite-energy endpoints, continuous approximations give
actual psh combinations $u_t=(1-t)u+tv$ by convexity of the
continuous cone and decreasing closure. For clarity, rational convex
combinations of FS metrics are FS by tensoring section systems after
clearing denominators. Uniform FS closure gives convexity of the
continuous cone, first for rational and then for real coefficients.
The continuous quadratic formula and positivity of $D$ give energy
concavity and a uniform finite lower bound for these combinations,
so their endpoints are in $\E^1$. Applied to the combined
approximations, Proposition \ref{prop:vector} gives
$[u_t]=(1-t)[u]+t[v]$. The test formula gives
$\mu_{u_t}=(1-t)\mu_u+t\mu_v$ by uniform density.
Finally \eqref{eq:singular-quadratic} is a quadratic polynomial
in $t$, whose derivative is
\[
 \int_T(v-u)\,d\mu_\psi-D([v]-[u],[u_t])
 =\int_T(v-u)\,d\mu_{u_t},
\]
using \eqref{eq:weighted} separately for $u$ and $v$.
Their finite integrals imply that both are finite on a common
full-measure set. The last expression therefore never subtracts
two infinities at a pole. The derivative includes one-sided
endpoint derivatives and ordinary interior derivatives. It is
an affine psh-segment derivative, not an outer envelope derivative
in every continuous direction.

\section{Reference changes and complete scalar extension}

If $\psi'$ is another continuous psh reference and
$r=\psi'-\psi$, the continuous energy cocycle, and then the
same absolute-majorant infimum, give
$E_{\psi'}(\phi)=E_\psi(\phi)-E_\psi(\psi')$.
Thus the actual finite-energy absolute metrics are independent of
the reference. A defining net of absolute continuous psh metrics
is the same under both references. Its continuous curvature measures
are the same, so Proposition \ref{prop:measure} gives the same
singular measure. The weighted conclusions concern these absolute
metrics and bounded reference changes; they therefore transport
by the ordinary integral law.

Let $F/K$ be any complete valued extension, and
$p:X_F^{\an}\to T$ the analytic projection. The curve remains
smooth and geometrically integral, and its geometric degree is $V$.
Continuous curvature satisfies $p_*\mu_{u_{i,F},c}=\mu_{u_i,c}$
and continuous energy is unchanged, by Proposition 8.3(iv),
Theorem 8.4 and Proposition 9.14(vi) of \cite{BE}.
The pulled-back defining net decreases to the actual pulled-back
potential. Lemma \ref{lem:energy} on both carriers proves its
extended energy equality, so $u_F\in\E^1_F$.
Apply Proposition \ref{prop:measure} to that net upstairs and
push its weak limit through the continuous projection. It gives
\[
 p_*\mu_{u_F}=\mu_u.
\]
No multiplication or division by the degree of a field extension
occurs. This is an ascent and pushforward statement, not descent
of every noninvariant upstairs potential. It completes the proof
of Theorem \ref{thm:main}.

\section{Exact remaining scope and source record}

The proof uses complete pertinent continuous primary units of
\cite{BE}: Definitions 5.15 and 7.2; Proposition 5.16, Lemma 5.17
and the whole Theorem 5.18 separation proof; Lemma 7.4 and
Corollaries 7.5--7.6; Theorem 7.14; the whole mixed-measure
Section 8.1 and its extension proof; Definition 9.13 and
Proposition 9.14. Smooth Stokes and positivity are the standard
CLD analytic foundations referenced by Proposition 8.3. No fresh
reading of every recursively cited foundation is claimed. A
separate literal-version source ledger records exact applications.
It is provenance, not independent mathematical certification.

No caller mathematical grants are used, and no existing
trivially-valued theorem is extended by changing its field name.
The direct ordered Dirichlet estimate is special to curves.
The argument proves neither an arbitrary-dimensional unbounded
mixed operator nor locality/comparison theorems or actual
existence and every-solution uniqueness for all original targets.
It supplies no normalized potential compactness, arbitrary envelope
continuity, orthogonality or both-sign outer envelope differentiation
over a nondiscrete field. Those original requirements remain open.
Even a finite-energy curvature operator alone is not a solution of
the whole original P16.4. This new draft must be reviewed in its
entirety before it enters the accepted paper count.

\begin{thebibliography}{9}
\bibitem{BE} S. Boucksom and D. Eriksson,
\emph{Norms, determinant of cohomology and Fekete points},
arXiv:1805.01016v3 (8 January 2021).
\url{https://arxiv.org/pdf/1805.01016v3}.
\end{thebibliography}
\end{document}
